\documentclass[
%% TIKZ_CLASSOPTION %%
tikz
]{standalone}
\usepackage{amsmath}
\usetikzlibrary{matrix}
%% EXTRA_TIKZ_PREAMBLE_CODE %%
\begin{document}
%% TIKZ_CODE %%
\begin{figure}
    \centering
    \tikzset{> = stealth,
    circ_targets/.style = {
        draw = black,
        fill=lightgray,
        shape = circle,
        inner sep = 1pt
    },
    circ/.style = {
        draw = black,
        shape = circle,
        inner sep = 1pt
    }}
    \tikz{
        \node[circ_targets] (D) at (0,0) {\bm{$D$}};
        \node[circ_targets] (y) at (2,0) {\bm{$Y$}};
        \node[circ] (x_tab) at (0,1) {$X_{\text{tab}}$};
        \node[circ] (x_txt) at (1,1.5) {$X_{\text{txt}}$};
        \node[circ] (x_img) at (2,1) {$X_{\text{img}}$};
        \node[circ] (U_tab) at (0,-1) {$U_{\text{tab}}$};
        \node[circ] (U_txt) at (1,-1.5) {$U_{\text{txt}}$};
        \node[circ] (U_img) at (2,-1) {$U_{\text{img}}$};
        \path[->, black] (x_tab) edge (D);
        \path[->, black] (x_tab) edge (y);
        \path[->, black] (U_tab) edge (D);
        \path[->, black] (U_tab) edge (y);
        \path[->, black] (x_txt) edge (D);
        \path[->, black] (x_txt) edge (y);
        \path[->, black] (U_txt) edge (D);
        \path[->, black] (U_txt) edge (y);
        \path[->, black] (x_img) edge (D);
        \path[->, black] (x_img) edge (y);
        \path[->, black] (U_img) edge (D);
        \path[->, black] (U_img) edge (y);
        \path[->, black, line width=1.5pt] (D) edge (y);
        \draw (D) -- (y) node [midway, fill=white] {$\theta_0$};
    }
    \caption{DAG for the semi-synthetic dataset. The confounding via the features $X=(X_{\text{tab}}, X_{\text{txt}}, X_{\text{img}})$ can be adjusted for, whereas the unexplained/noise parts $U=(U_{\text{tab}}, U_{\text{txt}}, U_{\text{img}})$ are unobserved.}
    \label{fig:conf_dag}
\end{figure}
\end{document}
